% kap05.tex — Kanoniske transformationer
% Hamilton-Jacobi-kompendiet

\chapter{Kanoniske transformationer}
\label{kap:kanoniske-transf}

\section{Motivation: frihed til at vælge faserumskoordinater}

I Lagrange-formalismen (kapitel~\ref{kap:lagrange}) kan vi vælge
generaliserede koordinater $q_i$ helt frit — enhver invertibel
omskrivning $Q_i=Q_i(q,t)$ giver stadig Euler-Lagrange-ligninger af
samme form. Spørgsmålet i dette kapitel er, om der findes en analog
frihed i Hamilton-formalismen: kan vi skifte til nye faserums\-
koordinater $(Q_i,P_i)$ — ikke kun nye $Q_i$, men også nye $P_i$, som
kan blande de gamle $q_i$ og $p_i$ på vilkårlig vis — og stadig have
Hamiltons ligninger på samme form,
\begin{equation}
  \dot Q_i = \pdv{K}{P_i}, \qquad \dot P_i = -\pdv{K}{Q_i},
  \label{eq:hamilton-nye-var}
\end{equation}
for en ny Hamilton-funktion $K(Q,P,t)$?

Svaret er ja, og mængden af sådanne transformationer er langt større
end i Lagrange-billedet — netop fordi $q_i$ og $p_i$ i Hamilton-
formalismen spiller en langt mere symmetrisk rolle end $q_i$ og
$\dot q_i$ gør i Lagrange-formalismen. Denne frihed er ikke blot en
teoretisk raffinement: det er selve kernen i strategien bag
Hamilton-Jacobi-teorien (kapitel~\ref{kap:hj}), hvor vi søger en
kanonisk transformation, der gør den \emph{nye} Hamilton-funktion
identisk nul — hvorved bevægelsesligningerne i de nye variable bliver
trivielle, $\dot Q_i=\dot P_i=0$.

\section{Definition}

En transformation $(q_i,p_i)\to(Q_i,P_i)$ (eventuelt tidsafhængig)
kaldes \emph{kanonisk}\index{kanonisk transformation}, hvis den bevarer Hamiltons ligningers form: der
findes en funktion $K(Q,P,t)$, så \eqref{eq:hamilton-nye-var} er
opfyldt for enhver Hamilton-funktion $H$, systemet kunne have. Vi skal
nu udlede en konstruktiv metode til at frembringe sådanne
transformationer — via \emph{genererende funktioner}.

\section{Det modificerede Hamiltons princip}

Hamiltons princip (afsnit~\ref{sec:hamiltons-princip}) kan omskrives i
faserumsvariable. Ved at bruge $L=\sum_i p_i\dot q_i - H$ (omskrivning
af definitionen \eqref{eq:def-H}), bliver virkningen
\begin{equation}
  S = \int_{t_1}^{t_2}\left(\sum_i p_i\dot q_i - H\right)dt,
\end{equation}
og variation af $S$ med hensyn til \emph{uafhængige} variationer i
$q_i$ og $p_i$ (ikke kun $q_i$, som i afsnit~\ref{sec:hamiltons-princip})
giver netop Hamiltons ligninger som Euler-Lagrange-ligninger for dette
udvidede variationsprincip (se opgave 1). For at de \emph{nye}
variable $(Q_i,P_i)$ skal opfylde Hamiltons ligninger med en ny
Hamilton-funktion $K$, skal det tilsvarende princip
\begin{equation}
  \delta\int_{t_1}^{t_2}\left(\sum_i P_i\dot Q_i - K\right)dt = 0
\end{equation}
gælde \emph{samtidig} med det oprindelige. Dette er sikret, hvis
integranderne kun adskiller sig ved en total tidsafledt af en funktion
$F$ — for et sådant led bidrager kun med randtermer, som ikke påvirker
variationen, når endepunkterne er fastholdt (præcis som i
udledningen af \eqref{eq:delta-S}):
\begin{equation}
  \sum_i p_i\dot q_i - H = \sum_i P_i\dot Q_i - K + \dv{F}{t}.
  \label{eq:genererende-betingelse}
\end{equation}
Funktionen $F$ kaldes den \emph{genererende funktion}\index{genererende funktion} for
transformationen. Den kan i princippet afhænge af en blanding af
gamle og nye variable — det er netop valget af, hvilke variable $F$
udtrykkes ved, der giver anledning til de fire ``typer'' af genererende
funktioner nedenfor.

\section{Genererende funktion af type $F_1(q,Q,t)$}

Vi vælger først $F=F_1(q,Q,t)$, en funktion af de \emph{gamle}
koordinater og de \emph{nye} koordinater. Da
\begin{equation}
  \dv{F_1}{t} = \sum_i\left(\pdv{F_1}{q_i}\dot q_i + \pdv{F_1}{Q_i}\dot Q_i\right) + \pdv{F_1}{t},
\end{equation}
bliver \eqref{eq:genererende-betingelse}
\begin{equation}
  \sum_i p_i\dot q_i - H = \sum_i P_i\dot Q_i - K
  + \sum_i\left(\pdv{F_1}{q_i}\dot q_i + \pdv{F_1}{Q_i}\dot Q_i\right) + \pdv{F_1}{t}.
\end{equation}
Da $q_i$ og $Q_i$ (og dermed $\dot q_i$, $\dot Q_i$) kan varieres
uafhængigt, må koefficienterne til $\dot q_i$ og $\dot Q_i$ stemme
overens hver for sig, og de resterende (ikke-tidsafledte) led må også
balancere:
\begin{equation}
  \boxed{p_i = \pdv{F_1}{q_i}, \qquad P_i = -\pdv{F_1}{Q_i}, \qquad K = H + \pdv{F_1}{t}.}
  \label{eq:F1-relationer}
\end{equation}
Givet en vilkårlig funktion $F_1(q,Q,t)$ definerer disse tre
relationer altså en kanonisk transformation: den første ligning
(løst for $Q_i$ som funktion af $(q,p,t)$) giver transformationen
$Q_i(q,p,t)$, den anden giver $P_i(q,p,t)$, og den tredje giver den nye
Hamilton-funktion.

\section{Genererende funktion af type $F_2(q,P,t)$}

I praksis er det ofte mere bekvemt at udtrykke den genererende
funktion ved de gamle koordinater og de \emph{nye impulser}, snarere
end de nye koordinater. Vi opnår dette ved endnu en Legendre-
transformation: definer
\begin{equation}
  F_2(q,P,t) \equiv F_1(q,Q,t) + \sum_i Q_i P_i,
\end{equation}
hvor $Q_i$ på højresiden forstås som funktion af $(q,P,t)$. Ved samme
type udregning som i kapitel~\ref{kap:hamilton} (differentier $F_2$
totalt, brug \eqref{eq:F1-relationer}, og sammenlign koefficienter —
se opgave 2 for detaljerne) finder man
\begin{equation}
  \boxed{p_i = \pdv{F_2}{q_i}, \qquad Q_i = \pdv{F_2}{P_i}, \qquad K = H + \pdv{F_2}{t}.}
  \label{eq:F2-relationer}
\end{equation}
Dette er den form, vi primært skal bruge i resten af kompendiet — og
i særdeleshed i kapitel~\ref{kap:hj}, hvor Hamilton-Jacobi-ligningen
netop er betingelsen for, at en genererende funktion af denne type
gør $K$ identisk nul.

\begin{bemaerkning}
De to øvrige typer, $F_3(p,Q,t)$ og $F_4(p,P,t)$, konstrueres ved
analoge Legendre-transformationer (med hensyn til henholdsvis $q$ og
$Q$) og giver relationerne
$q_i=-\partial F_3/\partial p_i$, $P_i=-\partial F_3/\partial Q_i$
og $q_i=-\partial F_4/\partial p_i$, $Q_i=\partial F_4/\partial P_i$
(begge med $K=H+\partial F/\partial t$). Vi vil ikke få brug for disse
i kompendiet, men de fire typer er blot fire forskellige ``udsnit'' af
samme genererende betingelse \eqref{eq:genererende-betingelse}, valgt
efter hvilket sæt variable der er mest bekvemt i en given situation.
\end{bemaerkning}

\section{Eksempel: identitetstransformationen og variableombytning}

\begin{eksempel}
Lad $F_2(q,P) = \sum_i q_iP_i$. Ifølge \eqref{eq:F2-relationer} er
\begin{equation}
  p_i = \pdv{F_2}{q_i} = P_i, \qquad Q_i = \pdv{F_2}{P_i} = q_i,
\end{equation}
altså den trivielle identitetstransformation. Dette bekræfter, at
konventionen i \eqref{eq:F2-relationer} er konsistent med den
oprindelige betegnelse af variablene.
\end{eksempel}

\begin{eksempel}
Lad i stedet $F_1(q,Q) = \sum_i q_iQ_i$ (bemærk: type $F_1$, ikke
$F_2$, denne gang). Ifølge \eqref{eq:F1-relationer} er
\begin{equation}
  p_i = \pdv{F_1}{q_i} = Q_i, \qquad P_i = -\pdv{F_1}{Q_i} = -q_i,
\end{equation}
altså transformationen $Q_i=p_i$, $P_i=-q_i$: de nye koordinater
\emph{er} de gamle impulser, og de nye impulser er (minus) de gamle
koordinater. Dette er et af de simplest mulige, men også mest
instruktive eksempler i hele kompendiet: det viser konkret, at
opdelingen i ``koordinater'' og ``impulser'' ikke er en fysisk nødvendighed,
men blot et bekvemt valg — Hamilton-formalismen behandler $q$ og $p$
symmetrisk, i modsætning til Lagrange-formalismen, hvor $q$ og
$\dot q$ har fundamentalt forskellige roller. (Bemærk minustegnet
på $P_i$: uden det ville $\{Q_i,P_j\}=\{p_i,-q_j\}=\delta_{ij}$ ikke
holde med korrekt fortegn — se opgave 3.)
\end{eksempel}

\section{Infinitesimale kanoniske transformationer og Poisson-parenteser}

Vi knytter nu en forbindelse tilbage til kapitel~\ref{kap:poisson}, som
bliver vigtig, når vi tænker på $H$ selv som ``generator'' af
tidsudvikling. Betragt en genererende funktion af formen
\begin{equation}
  F_2(q,P) = \sum_i q_iP_i + \varepsilon\,G(q,P),
\end{equation}
hvor $\varepsilon$ er en infinitesimal parameter, og $G$ er en
vilkårlig (glat) funktion — en lille perturbation af
identitetstransformationen. Ifølge \eqref{eq:F2-relationer}, til
første orden i $\varepsilon$:
\begin{equation}
  p_i = P_i + \varepsilon\pdv{G}{q_i} \quad\Longrightarrow\quad P_i = p_i - \varepsilon\pdv{G}{q_i} + O(\varepsilon^2),
\end{equation}
\begin{equation}
  Q_i = q_i + \varepsilon\pdv{G}{P_i} = q_i + \varepsilon\pdv{G}{p_i} + O(\varepsilon^2)
\end{equation}
(i sidste skridt har vi til denne orden erstattet $P_i$ med $p_i$
inde i $\partial G/\partial P_i$, hvad der kun ændrer resultatet med et
$O(\varepsilon^2)$-bidrag). Ændringen i en vilkårlig funktion
$f(q,p)$ under denne infinitesimale transformation er derfor
\begin{equation}
  \delta f \equiv f(Q,P) - f(q,p)
  = \sum_i\left(\pdv{f}{q_i}\,\varepsilon\pdv{G}{p_i} - \pdv{f}{p_i}\,\varepsilon\pdv{G}{q_i}\right) + O(\varepsilon^2)
  = \varepsilon\,\{f,G\}.
  \label{eq:infinitesimal-generator}
\end{equation}
Vi genkender netop Poisson-parentesen fra kapitel~\ref{kap:poisson}!
Funktionen $G$ kaldes \emph{generatoren}\index{generator (infinitesimal kanonisk transformation)} af transformationen, og
\eqref{eq:infinitesimal-generator} viser, at Poisson-parentesen
$\{f,G\}$ måler, hvordan $f$ ændrer sig under den infinitesimale
kanoniske transformation genereret af $G$.

\begin{bemaerkning}
Dette giver en dybere fortolkning af flere resultater, vi allerede har
set. I opgave 3 til kapitel~\ref{kap:poisson} fandt vi
$\{L_z,x\}=-y$, $\{L_z,y\}=x$ — netop hastigheden, hvormed $(x,y)$
ændrer sig under en infinitesimal rotation om $z$-aksen. Med
\eqref{eq:infinitesimal-generator} ser vi, at dette \emph{er} den
generelle udsagn: $L_z$ genererer infinitesimale rotationer om
$z$-aksen som en kanonisk transformation, med $\{f,L_z\}$ som ``hastigheden'',
hvormed en given størrelse $f$ ændrer sig under rotationen.

Endnu vigtigere for kompendiets videre gang: sæt $G=H$ og
$\varepsilon = dt$. Da bliver \eqref{eq:infinitesimal-generator}
$\delta f = dt\,\{f,H\}$, altså netop bevægelsesligningen
\eqref{eq:bevaegelse-poisson} (for tidsuafhængige $f$). \emph{Selve
tidsudviklingen er en kontinuert kanonisk transformation, genereret af
Hamilton-funktionen selv.} Dette er den konceptuelle kerne, vi bygger
videre på i kapitel~\ref{kap:hj}: hvis vi kan finde en genererende
funktion, der udfører \emph{hele} tidsudviklingen på én gang (fra
tid $t_1$ til $t_2$) som én kanonisk transformation, vil den nye
Hamilton-funktion være identisk nul, og den nye ``position'' $Q_i$ og
``impuls'' $P_i$ vil simpelthen være konstanterne, de startede med. Den
funktion er netop Hamiltons principfunktion $S$.
\end{bemaerkning}

\section{Kanoniske transformationer bevarer Poisson-parenteserne}

Vi angiver, uden fuldt bevis (se opgave 4 for et specialtilfælde), et
centralt strukturelt resultat: en transformation $(q,p)\to(Q,P)$ er
kanonisk, hvis og kun hvis den bevarer de fundamentale Poisson-
parenteser,
\begin{equation}
  \{Q_i,Q_j\}_{q,p} = 0, \qquad \{P_i,P_j\}_{q,p} = 0, \qquad \{Q_i,P_j\}_{q,p} = \delta_{ij},
  \label{eq:kanonisk-poisson-betingelse}
\end{equation}
hvor $\{\cdot,\cdot\}_{q,p}$ betyder, at parentesen \eqref{eq:def-poisson}
udregnes med afledte med hensyn til de \emph{gamle} variable $(q,p)$.
Mere generelt bevarer en kanonisk transformation \emph{alle} Poisson-
parenteser: for to funktioner $f,g$, udtrykt enten ved $(q,p)$ eller
ved de transformerede $(Q,P)$, er $\{f,g\}_{q,p}=\{f,g\}_{Q,P}$. Dette
er den dybere grund til, at Poisson-parentesen er så nyttig et
redskab: den er ikke blot en bekvem notation i ét bestemt
koordinatsystem, men en \emph{invariant} struktur på faserummet,
uafhængig af hvilke kanoniske koordinater vi måler den i.

\begin{bemaerkning}
Denne invarians er faserummets analog til, hvordan volumenet i
Liouvilles sætning (kapitel~\ref{kap:poisson}) er bevaret — begge
udsagn siger, at kanoniske transformationer respekterer en underliggende
geometrisk struktur (den \emph{symplektiske} struktur) på faserummet,
uanset hvor ``vildt'' transformationen blander $q$'er og $p$'er. Det er
denne bevarede struktur, der gør faserummet — og ikke bare et vilkårligt
$2n$-dimensionalt rum af tal — til den rigtige arena for mekanik.
\end{bemaerkning}

\begin{figure}[htbp]
  \centering
  \begin{tikzpicture}[scale=1.0]
    \draw[-{Latex}] (-0.3,0) -- (3.4,0) node[right] {$q$};
    \draw[-{Latex}] (0,-0.3) -- (0,3.0) node[above] {$p$};
    \draw[thick,blue!60!black,fill=blue!12] (0.8,0.8) rectangle (1.8,1.8);
    \node[blue!60!black] at (1.3,1.3) {$\Omega$};
    \node at (1.3,-0.7) {$(q,p)$};

    \draw[-{Latex},thick] (4.0,1.3) -- (5.4,1.3) node[midway,above] {kanonisk};

    \draw[-{Latex}] (6.0,0) -- (9.7,0) node[right] {$Q$};
    \draw[-{Latex}] (6.0,-0.3) -- (6.0,3.0) node[above] {$P$};
    \draw[thick,black,fill=black!8] plot[smooth cycle] coordinates
      {(7.0,0.9) (8.1,0.7) (8.7,1.5) (8.0,2.2) (7.1,1.7)};
    \node at (7.9,1.35) {$\Omega'$};
    \node at (7.9,-0.7) {$(Q,P)$};
  \end{tikzpicture}
  \caption{En kanonisk transformation kan forvrænge et område $\Omega$ i
  faserummet ganske voldsomt i form — men arealet (i to dimensioner)
  hhv.\ volumenet (i $2n$ dimensioner) er altid bevaret, ligesom under
  den rene tidsudvikling i Liouvilles sætning
  (figur~\ref{fig:liouville}).}
  \label{fig:kanonisk-transf}
\end{figure}

\section{Opsummering}

\begin{itemize}
  \item En kanonisk transformation $(q,p)\to(Q,P)$ er en transformation,
        der bevarer Hamiltons ligningers form, med en ny
        Hamilton-funktion $K$.
  \item Sådanne transformationer frembringes af genererende funktioner
        $F$, defineret via
        $\sum_ip_i\dot q_i-H=\sum_iP_i\dot Q_i-K+dF/dt$.
  \item For $F_1(q,Q,t)$: $p_i=\partial F_1/\partial q_i$,
        $P_i=-\partial F_1/\partial Q_i$, $K=H+\partial F_1/\partial t$.
  \item For $F_2(q,P,t)=F_1+\sum_iQ_iP_i$: $p_i=\partial F_2/\partial q_i$,
        $Q_i=\partial F_2/\partial P_i$, $K=H+\partial F_2/\partial t$
        — den type, vi primært bruger fremover.
  \item En infinitesimal kanonisk transformation genereret af $G$
        giver $\delta f=\varepsilon\{f,G\}$; specielt genererer $H$
        selv (med $\varepsilon=dt$) den faktiske tidsudvikling.
  \item Kanoniske transformationer bevarer alle Poisson-parenteser —
        Poisson-strukturen er en invariant egenskab af faserummet, ikke
        af et bestemt koordinatvalg.
\end{itemize}

\begin{opgave}
Vis, at variation af $S=\int(\sum_ip_i\dot q_i-H)\,dt$ med hensyn til
uafhængige, arbitrære $\delta q_i(t)$ og $\delta p_i(t)$ (begge nul
ved endepunkterne $t_1,t_2$ for $\delta q_i$; ingen randbetingelse er
nødvendig for $\delta p_i$, da der ikke integreres partielt i $p_i$)
giver netop Hamiltons ligninger \eqref{eq:hamiltons-ligninger} som
Euler-Lagrange-ligninger.
\end{opgave}

\begin{opgave}
Udfør mellemregningen, der fører fra $F_1$-relationerne
\eqref{eq:F1-relationer} til $F_2$-relationerne
\eqref{eq:F2-relationer}, ved at differentiere
$F_2=F_1+\sum_iQ_iP_i$ totalt og sammenligne med
\eqref{eq:genererende-betingelse} skrevet om for $F_2$.
\end{opgave}

\begin{opgave}
Bekræft for eksemplet $Q_i=p_i$, $P_i=-q_i$, at de fundamentale
Poisson-parenteser \eqref{eq:kanonisk-poisson-betingelse} er opfyldt
(udregnet med hensyn til de gamle variable $(q,p)$), og vis, at
transformationen \emph{uden} minustegnet (dvs.\ $Q_i=p_i,P_i=q_i$)
ville bryde $\{Q_i,P_i\}=\delta_{ij}$.
\end{opgave}

\begin{opgave}
Lad $G=H$ selv være generator af en infinitesimal transformation med
$\varepsilon=dt$, og brug \eqref{eq:infinitesimal-generator} med
$f=q_i$ og $f=p_i$ til at genfinde Hamiltons ligninger. Diskuter,
hvorfor dette bekræfter påstanden i bemærkningen om, at ``tidsudvikling
er en kanonisk transformation genereret af $H$''.
\end{opgave}
